@@M0@@ the augmentation homomorphism: @@M1@@.

\sgaStableTarget{sga6.e06.prop.6-1}\textbf{Proposition 6.1.} \uline{Let} @@M2@@ \uline{be a quasi}-\uline{compact scheme having an ample sheaf} @@M3@@. \uline{Then} @@M4@@ \uline{is a nil ideal}.

Since @@M5@@ is generated by elements of the form @@M6@@, where @@M7@@ is the class of an @@M8@@-Module locally free of finite type @@M9@@, and @@M10@@ is the augmentation of @@M11@@, it suffices to show that such an element is nilpotent. The quasi-compactness of @@M12@@, and the commutation of @@M13@@ with finite direct sums, reduce us to the case where @@M14@@ has constant rank. Finally, passing to the flag bundle of @@M15@@, which gives an injection on the @@M16@@ (3.1) and preserves the hypothesis of the existence of an ample sheaf, we may suppose that @@M17@@ is invertible.

Let @@M18@@ be the dual of @@M19@@. For @@M20@@ sufficiently large, @@M21@@ is generated by its global sections; we can therefore find an integer @@M22@@ and a surjective homomorphism: @@M23@@, hence (tensoring by @@M24@@) a surjective homomorphism: @@M25@@. The kernel of this homomorphism is locally free of rank @@M26@@. Thus, denoting by @@M27@@ the class of @@M28@@, we obtain:
@@M29@@
Now @@M30@@. Consequently, for every invertible sheaf @@M31@@, the element @@M32@@ is nilpotent for @@M33@@ sufficiently large, or equivalently its product @@M34@@ by @@M35@@ is nilpotent.

We can then write:
@@M36@@
and, by what precedes, @@M37@@ and @@M38@@ are nilpotent on choosing @@M39@@ sufficiently large. Thus @@M40@@ is nilpotent.

